\chapter{Literature Review}
\label{ch:literature}

\section{Introduction}
\label{sec:lit-intro}

This chapter reviews the literature underpinning this project in four parts. Section~\ref{sec:lit-fl} introduces federated learning and the poisoning attack surface it creates. Section~\ref{sec:lit-defenses} reviews Byzantine-robust and statistical defenses, the primary line of defense this report benchmarks against. Section~\ref{sec:lit-attacks} reviews adaptive and learning-based poisoning attacks, the category of attack this project extends. Section~\ref{sec:lit-rl} reviews the reinforcement-learning and parameter-efficient fine-tuning methods used to train the LLM policies in this report. Section~\ref{sec:lit-gap} closes the chapter by identifying the specific gap this project addresses.

\section{Federated Learning and the Poisoning Attack Surface}
\label{sec:lit-fl}

In a standard FL round, a central server distributes the global model to $N$ clients, each of which performs local training on its own data and returns an update, which the server aggregates -- typically through FedAvg -- into a new global model \citep{fedavg}. Because the server never inspects a client's raw data, model poisoning is possible: a compromised client (or a small set of them) submits a manipulated update instead of an honestly trained one. \emph{Untargeted} poisoning attacks aim to degrade the global model's overall accuracy; \emph{targeted}, or backdoor, attacks instead aim to introduce a narrow vulnerability -- misclassifying one class, or inputs containing a specific trigger pattern -- while leaving overall accuracy largely intact, so the attack is harder to notice from aggregate metrics alone \citep{bagdasaryan2020}. This report focuses on the untargeted case.

Non-IID data partitioning -- where each client's local data distribution differs from the population distribution, which is the norm rather than the exception in real FL deployments -- substantially complicates defense. Under a non-IID partition, an honest client's update can legitimately resemble a statistical outlier simply because its local class distribution departs from the majority, which is exactly the signal many defenses use to identify malicious clients. Any defense (or attack) evaluated only under an IID partition risks over-stating its real-world effectiveness; for this reason, this report uses a non-IID partition throughout (Section~\ref{sec:method-environment}).

\section{Byzantine-Robust and Statistical Defenses}
\label{sec:lit-defenses}

\subsection{Coordinate- and Distance-Based Methods}
Coordinate-wise median and trimmed-mean aggregation provide provable robustness guarantees under a bounded number of Byzantine clients by limiting the influence any single coordinate of the update vector can have on the aggregate \citep{byzantine}. Krum, and its multi-client extension Multi-Krum, take a different, distance-based approach: each client is scored by the sum of squared Euclidean distances to its closest neighbours, and only the clients with the smallest (most "central") scores are kept and averaged \citep{krum}. This makes Multi-Krum effective against updates that are geometric outliers relative to the honest majority, but, as later sections of this report show empirically, it also depends on being given (or assuming) a reasonably accurate upper bound on the number of malicious clients.

\subsection{Spectral Methods}
DnC (Divide-and-Conquer) treats poisoning as a spectral outlier problem rather than a purely geometric one: it centres the client updates, computes the top right singular vector of the centred update matrix via singular value decomposition, projects every client's centred update onto that direction, and removes the clients with the largest squared projection \citep{dnc}. Because the direction a model-poisoning attacker is often forced to move the aggregate along tends to dominate the top singular direction of the update matrix, this can catch attacks that a purely pairwise-distance method might miss.

\subsection{Trust-Bootstrapped Methods}
FLTrust departs from the fully unsupervised setting of the methods above by allowing the server to hold a small, clean root dataset \citep{fltrust}. Each round, the server fine-tunes the current global model on this root set to obtain a trusted reference update, scores every client update by the rectified cosine similarity to that reference, rescales surviving updates to the reference update's magnitude (which specifically neutralises attacks that simply scale up a malicious update to dominate the average), and combines the rescaled updates in a trust-weighted average. Because the trust score is continuous and the combination step is a weighted average rather than a hard accept/reject filter, FLTrust's real robustness does not depend on making a sharp binary classification decision -- a property this report returns to directly in its results (Section~\ref{sec:results-detection-vs-protection}).

\subsection{Similarity-Based Methods}
FoolsGold takes the opposite approach to the methods above: rather than looking for outliers, it looks for clients whose update \emph{histories} are unusually similar to one another, on the assumption that colluding Sybil clients will tend to move in similar directions round after round, and down-weights clients accordingly \citep{foolsgold}. This is the basis for the collusion-aware terms used elsewhere in this report's reward design (Section~\ref{sec:method-rewards}).

\subsection{Layer-Profile-Based Methods}
DeFL differs from every method above in that it does not compute a single per-round, whole-update statistic at all. Instead, it tracks a per-layer Federated Gradient Norm Vector across rounds, detects a \emph{Critical Learning Period} from a sufficiently large round-over-round rise in the aggregate per-layer gradient norm, and, during such a period, votes a client malicious using an adaptively falling per-layer outlier threshold \citep{defl}. This makes DeFL's filtering strength time-varying by design -- a property this report's results (Section~\ref{sec:results-detection-vs-protection}) suggest may leave it vulnerable to damage that compounds across rounds outside a detected critical period, even when its per-round detection accuracy looks strong.

\section{Adaptive and Learning-Based Attacks}
\label{sec:lit-attacks}

A separate body of work shows that an attacker aware of a specific deployed defense can construct an update that satisfies that defense's statistical assumptions while still corrupting the aggregate. Fang~et~al.\ formulate local model poisoning directly as an optimisation problem against a known Byzantine-robust aggregation rule, showing that several rules assumed robust can still be substantially degraded by an attacker who solves for the worst-case update permitted under that rule \citep{fang2020}. Bhagoji~et~al.\ show that even a single, persistent attacker using a boosted update can meaningfully bias the global model over many rounds, even under conventional defenses \citep{bhagoji2019}. Bagdasaryan~et~al.\ demonstrate that a backdoor can be implanted into a federated model while the model's overall accuracy remains high, illustrating how narrow the margin can be between an attack a defense will and will not catch \citep{bagdasaryan2020}.

These attacks share a common limitation for the purposes of this project: each is hand-derived against one specific, known defense. This report instead trains an attacker through interaction and a verifiable reward signal, with no white-box access to any defense's parameters or thresholds -- a methodological difference discussed further in Section~\ref{sec:lit-gap}.

\section{Reinforcement Learning for LLM Policies}
\label{sec:lit-rl}

\subsection{Policy-Gradient Methods}
Proximal Policy Optimization (PPO) is a widely used policy-gradient algorithm for aligning and specialising LLMs, and trains a learned value (critic) network alongside the policy to provide per-token credit assignment along a generated sequence \citep{ppo,ouyang2022}. Direct Preference Optimization (DPO) instead learns directly from pairwise preference data, without an explicit reward model or online rollouts \citep{dpo}; this makes it simple and stable, but it is an offline method and discards the \emph{magnitude} of a reward signal, retaining only its ordering. Group-Relative Policy Optimization (GRPO) offers a third alternative: it samples several completions for the same input, scores each with an exactly computable (verifiable) reward, and uses the reward's deviation from the group's own mean as the training signal -- removing the need for a separate learned critic network altogether \citep{grpo}. This critic-free, online, magnitude-aware combination of properties is why GRPO was selected for this project (Section~\ref{sec:method-rl}): each round in the FL environment yields exactly one terminal, exactly verifiable reward per agent, which is precisely the setting GRPO was designed for.

\subsection{Parameter-Efficient Fine-Tuning}
Low-Rank Adaptation (LoRA) freezes a pretrained model's weights and injects small, trainable low-rank update matrices into selected layers, dramatically reducing the number of trainable parameters needed to specialise a large model \citep{lora2022}. Its 4-bit quantised variant, QLoRA, further reduces the memory footprint required to fine-tune a given base model \citep{qlora2023}. Both techniques make it practical for this project's two policies -- the attacker and the defender -- to share a single frozen base LLM behind two small, independently trained adapters, rather than requiring two full copies of the base model (Section~\ref{sec:method-impl}).

\section{Research Gap}
\label{sec:lit-gap}

Two gaps, taken together, motivate this report. First, prior adaptive attacks \citep{fang2020,bhagoji2019} are hand-derived against one specific, known defense and its parameters; none of the reviewed literature trains an attack policy through interaction and a verifiable reward, with no white-box access to the defense at all, and then tests whether the resulting strategy transfers to defenses it never saw. Second, the classical defenses reviewed in Section~\ref{sec:lit-defenses} are typically evaluated against different, non-uniform attacks across different papers, which makes their relative robustness difficult to compare directly; a single learned attack, replayed unmodified against several defenses under identical conditions, would make that comparison direct rather than indirect. This report addresses both gaps together: it trains an attacker with reinforcement learning against a co-trained defender it can only observe through aggregate outcomes (Chapter~\ref{ch:methodology}), and then benchmarks the resulting, frozen attacker against a panel of classical and adaptive defenses under one fixed-attack, varied-defense protocol (Chapter~\ref{ch:results}).
